\section{Scope of the Review}

This review covers four things: the benchmark this thesis uses, the architectures it
compares, the class-imbalance remedies it crosses them with, and the statistical method
it uses to decide whether two models actually differ. Work on segmentation, multimodal
fusion and explainability is noted where it bears on the comparison but is not the
subject here.

\section{HAM10000 as a Benchmark}
\label{sec:lit-ham}

HAM10000 was released by Tschandl et al.~\cite{tschandl2018ham10000} as a deliberate
answer to the shortage of public dermatoscopic data. It contains 10{,}015 images across
seven diagnostic categories, collected over two decades from a Vienna clinic and an
Australian practice, and it became the training set for the ISIC 2018
challenge~\cite{codella2019isic}. Ground truth is more defensible than in most
comparable sets: over half the cases are histopathologically confirmed rather than
assigned by consensus.

Two properties matter for what follows. The first is that images are indexed by
\texttt{lesion\_id} as well as \texttt{image\_id}, because the same lesion is often
photographed more than once. Any split that ignores this puts near-duplicates on both
sides of the train/test boundary. The second is the class distribution, which is
severely long-tailed: melanocytic nevi account for roughly two thirds of the set while
dermatofibroma accounts for about one percent. The ratio between the most and least
frequent class is close to 58 to 1.

The literature is inconsistent about both. Reported HAM10000 accuracies range widely,
and a meaningful part of that spread comes from protocol rather than method. Some studies
split at image level, some at lesion level, and many do not say. Gao et
al.~\cite{gao2024skinlitenet} report 98.97\% accuracy with a 284K-parameter network,
which is difficult to reconcile with the difficulty other work reports on the same data
unless the evaluation protocols differ substantially. This thesis takes the view that
cross-paper accuracy comparison on HAM10000 is largely uninformative, and that the only
comparisons worth trusting are ones run inside a single controlled pipeline.

\section{Architectures}
\label{sec:lit-arch}

\subsection{Residual networks as the default baseline}

He et al.~\cite{he2016resnet} introduced residual connections to make very deep networks
trainable, and ResNet-50 has been the default medical-imaging backbone since. In
dermatology specifically it appears constantly: Rezvantalab et
al.~\cite{rezvantalab2020dermoscopy} benchmark it among four architectures on
dermoscopic data, and Shams et al.~\cite{shams2025jcmr} use it as the reference point
against lighter alternatives. Its ubiquity is the reason it serves as the heavyweight
arm here. It is what a practitioner would reach for by default, so it is the right thing
for a lightweight candidate to be measured against.

\subsection{Efficiency-oriented designs}

Howard et al.~\cite{howard2017mobilenets} introduced depthwise separable convolutions to
cut the cost of a convolution by roughly an order of magnitude, and Sandler et
al.~\cite{sandler2018mobilenetv2} refined this into MobileNetV2 with inverted residuals
and linear bottlenecks. The design target was explicitly on-device inference.

The dermatology literature has taken this up. Hussain et al.~\cite{eswa2024e2e} propose
S-MobileNet for seven-class HAM10000 classification, and Gao et
al.~\cite{gao2024skinlitenet} push further with a 284K-parameter module. Both report
strong numbers. Neither reports a controlled comparison against a heavyweight baseline
under identical conditions, which is the specific thing this thesis is built to provide.

\subsection{EfficientNet and compound scaling}

EfficientNet occupies a middle position: more capacity than MobileNetV2, far less than
ResNet-50, with depth, width and input resolution scaled together rather than
independently. Manole et al.~\cite{manole2024efficientnet} apply it to dermatological
diagnosis and report 95.4\% validation accuracy on a four-class problem. Shams et
al.~\cite{shams2025jcmr} compare ResNet-50, MobileNet and EfficientNet-B0 on DermNet and
report 93\%, 94\% and 99\% respectively, which would make EfficientNet the clear choice
if the comparison were controlled.

It is not, and that is precisely the difficulty. EfficientNet variants are trained at
their own native input resolutions, which increase with the variant index. B3 expects
300 pixels where MobileNetV2 and ResNet-50 expect 224. A study that gives EfficientNet
its native resolution and the baselines theirs has varied two things at once, and cannot
separate architectural merit from input size. This thesis addresses that directly by
running B3 at both resolutions with everything else held fixed, which is what makes the
pair E7/E10 informative rather than merely additional.

\section{Class Imbalance}
\label{sec:lit-imbalance}

Buda et al.~\cite{buda2018classimbalance} provide the most systematic treatment,
examining imbalance across MNIST, CIFAR-10 and ImageNet. Three findings carry over.
Imbalance measurably degrades CNN performance; oversampling is generally the strongest
of the standard remedies; and oversampling does not induce the overfitting in CNNs that
it can in classical models. The last point is the one most often assumed rather than
tested, and it is worth testing on a dataset as skewed as HAM10000.

Cui et al.~\cite{cui2019classbalanced} reweight by effective number of samples,
generalising inverse-frequency weighting. That family of methods is what the weighted
cross-entropy condition in this thesis implements, using $w_c = N / (K \cdot n_c)$. On
the split used here that produces weights spanning 0.214 for nevi to 12.578 for
dermatofibroma, a spread of roughly 59 to 1. Weights of that magnitude are worth
treating as an intervention with side effects rather than a neutral correction.

Lin et al.~\cite{lin2017focal} propose focal loss, down-weighting well-classified
examples so training concentrates on hard ones. Le et al.~\cite{le2020classweightedfocal}
combine it with class weighting on HAM10000 and report 93\% top-1 accuracy averaged
across the seven classes. Focal loss is deliberately out of scope here, for reasons of
factorial size rather than merit: adding it as a fourth strategy would take the design
from twelve runs to sixteen. It is cited because a reviewer will ask, and because Le et
al.\ is the closest published comparator to this work.

\section{The Gap Le et al.\ Leaves Open}
\label{sec:lit-gap}

Le et al.~\cite{le2020classweightedfocal} is the nearest neighbour to this thesis and
therefore the most useful thing to be precise about. They vary the loss function on
HAM10000 while holding the architecture fixed at a modified ResNet-50, and they report
accuracy averaged over classes.

That leaves two things unanswered. The architecture never varies, so the study says
nothing about what a lighter backbone would cost. And averaged accuracy hides the
per-class picture, which is where the clinically relevant behaviour lives. A model can
improve its average while getting worse at melanoma, and nothing in an averaged figure
would reveal it.

The same pattern holds more broadly. Across the surveyed corpus, studies vary one factor
at a time (architecture, or loss, or sampling) and report an aggregate metric. No
published HAM10000 result crosses architecture with imbalance strategy in a single
controlled design while reporting per-class results for the malignant categories.

\section{Deciding Whether Two Models Differ}
\label{sec:lit-stats}

Comparing two classifiers on one test set is a statistical question, and the wrong test
is common. Dietterich~\cite{dietterich1998tests} evaluated five approximate tests for
this purpose and found that when each model is trained once and evaluated on a single
held-out set, McNemar's test is the only one with acceptable Type-I error. Tests
assuming independent samples do not apply, because both models are evaluated on the same
images.

That regime is exactly this thesis's. Each architecture-strategy combination is one
training run producing one set of predictions. McNemar's test operates on the discordant
pairs, the images one model classifies correctly and the other does not, and ignores the
cases where both agree. Agreement carries no information about which model is better.

Running many such comparisons inflates the family-wise error rate, so correction is
required. Holm-Bonferroni is used here as a uniformly more powerful alternative to plain
Bonferroni that makes no additional assumptions about dependence between tests.

\section{What the Review Establishes}

Four things follow from the above and shape the rest of this thesis. Cross-paper
accuracy comparison on HAM10000 is unreliable, so the comparison must happen inside one
pipeline. Per-class reporting is necessary because aggregate metrics hide the clinically
important behaviour. Input resolution is a confound whenever EfficientNet is compared
against 224-pixel baselines, and must be controlled explicitly. And McNemar's test with
correction across a declared family is the appropriate way to decide whether an observed
difference is real.
